\subsection{The harmonic sign is unchanged by cubic descent}

\paragraph{D23.P13 (an actual-source relation between the two descents).}
For every actual multiplier-one gentle element $h=(1,f)$,
the cubic descent of D23.F22 satisfies
\[
 \boxed{\quad a_2(\mathcal D_Eh)=a_2(h),\qquad
 \rho(\mathcal D_Eh)=\rho(h)\pmod3.\quad}
\tag{D23.P13.1}
\]
More precisely its normalized pair coordinate $f_E$ has exactly
the same value as $f$ in the signed affine harmonic detector of
supplied FP.34. Consequently, using D23.F22 and supplied A.107,
\[
 \boxed{\quad
 \Delta\circ\mathcal D_E
 =a_2+\Psi_3/2
 =d_2
 =a_2+a_2\circ D_L
 =a_2+a_2\circ D_P
 \quad\text{on }\mathcal H.\quad}
\tag{D23.P13.2}
\]
The half-character in this display is reduced modulo two.
This identifies an actual descended defect; it does not assert
that $d_2$ vanishes, or that either descent fixes all full-GT
elements.

\begin{proof}
We keep the exact FP.34 marking. Put
\[
 R=(\mathbb Z/4)[\omega]/(\omega^2+\omega+1),\qquad
 c=(1-\omega)^{-1}=2+3\omega,
\]
and use the affine multiplication $(a,u)(b,v)=(a+ub,uv)$.
The homomorphism $B=C_2\amalg C_3\to R\rtimes\langle-\omega\rangle$
is
\[
 s\longmapsto(c,-1),\quad r\longmapsto(c,\omega),\quad
 x\longmapsto(0,\omega),\quad y\longmapsto(1,\omega).
\tag{D23.P13.3}
\]
Write $f((0,\omega),(1,\omega))=(k,1)$.

For any commutator word $f$ and any $a,b\in R$, the exact
affine identity
\[
 f((a,\omega),(b,\omega))=((b-a)k,1)
\tag{D23.P13.4}
\]
holds. Indeed simultaneous conjugation by the translation
$ac$ changes $(0,\omega),(b-a,\omega)$ to the indicated
pair. A commutator word has linear part one, and conjugation
by a translation leaves a pure translation unchanged.
Finally multiplying the translation module by $b-a$ is an
affine-group homomorphism (no invertibility of $b-a$ is needed),
which gives (D23.P13.4). The argument applies to profinite
words by continuity, since the target is finite.

Use the notation and normalized formula from D23.F22:
\[
 C=f(x,r),\qquad g=f(r,srs),\qquad t=3\Psi_3(f)/2,
 \qquad f_E=y^{-t}(sC^{-1}s)gCx^t.
\tag{D23.P13.5}
\]
Here $C,g\in F$ and $t\in3\widehat{\mathbb Z}$; the latter
uses the already proved evenness of $\Psi_3$.
Since $srs\mapsto(-\omega c,\omega)$, formula (D23.P13.4) gives
\[
 C\longmapsto(ck,1),\qquad
 g\longmapsto(\omega^2ck,1).
\tag{D23.P13.6}
\]
Conjugation by $s$ negates translations, hence
$sC^{-1}s\mapsto(ck,1)$. Therefore
\[
 (sC^{-1}s)gC\longmapsto((2+\omega^2)ck,1)=(k,1),
\]
because $2+\omega^2=1-\omega$ and $(1-\omega)c=1$ in $R$.
Both displayed images of $x,y$ have order three. The normalizing
powers $x^t,y^{-t}$ thus have identity image, and
\[
                    q_{\rm harm}(f_E)=q_{\rm harm}(f).
\tag{D23.P13.7}
\]
This is an exact equality of finite values, not just a comparison
after forgetting a sign or reducing modulo two.

FP.34 gives, for each actual gentle normalized pair,
\[
 1-(1-\omega)k=\omega^\rho(-1)^{a_2}.
\]
The six elements $\omega^j(-1)^e$, $j\in\mathbb F_3$,
$e\in\mathbb F_2$, are distinct in $R$. Applying the formula
to $f_E$ and $f$ proves both claims in (D23.P13.1), independently
of the second congruence proved by the cubic matrices in D23.F22.

D23.F22 gives $a_3\circ\mathcal D_E=\Psi_3/2$. Adding
(D23.P13.1) proves the first two equalities of (D23.P13.2).
The remaining equalities are exactly supplied A.107, equations
(MD.8a)--(MD.8b). This deduction uses genuine endomorphisms
and actual characters on $\mathcal H$ throughout; it makes
no assertion about full-GT descent compatibility.

The exact certificate \path{agents/pentagon/cubic_descent_harmonic.py}
checks the ring, all signed peripheral images, the affine group
order, the translation computation for every $k\in R$, and
the uniqueness of the six harmonic phases. It does not replace
the cited actual-source descent construction in D23.F22 or the
supplied harmonic-frame theorem FP.34.
\end{proof}

\paragraph{Exact source locators and remaining question.}
The signed harmonic representation and recovery formula are
\path{sources/DELTA_ON_GT/agent_five/five_point_results.tex},
FP.34.1--FP.34.3. The pair formula is
\path{work/agents/five_level_six/cubic_descent.tex},
D23.F22.1 and its displayed normalization of $f_E$.
The half-character comparison with $D_L,D_P$ is
\path{sources/DELTA_ON_GT/sections/A.tex}, A.107,
(MD.8a)--(MD.8b).
The identity $\Psi_3/2\circ\mathcal D_E=\Psi_3/2$, or a different
formula for this iterated half-character, is not proved here.
